% GENERATED FILE. Edit canonical lessons, then run npm run generate:books.
% canonical-source-hash: fnv1a64:23231554f64c735e
% canonical-lessons: GE-C518-gemaelde, GE-C518-pinsel, GE-C518-leinwand, GE-C518-staffelei, GE-C518-aquarell

\chapter{Pinsel und Leinwand}
\label{ch:ge-pinsel-und-leinwand}
\hlchaptermodality{518}

\begin{chapteropening}
\textbf{By the end of this chapter:} \emph{I can talk about painting.}

Complete the last lesson of chapter 518: I can talk about painting.
\end{chapteropening}

\section[das Gemälde]{das Gemälde — painting}

\label{lesson:GE-C518-gemaelde}



\begin{quote}
Before the new one: say the German for to applaud, then the German for applause.
\end{quote}

\subsection*{You'll want to know: das Gemälde}
\textbf{das Gemälde} — \textquotedblleft{}painting\textquotedblright{}.

\begin{quote}
Im Museum hängt ein berühmtes Gemälde aus dem Mittelalter.
\end{quote}

\begin{quote}
\emph{A famous painting from the Middle Ages hangs in the museum.}
\end{quote}

Plural: \emph{die Gemälde}. Built on \emph{malen} — to paint.

\begin{cousinweb}
\emph{Gemälde} is built on \emph{malen} \textquotedblleft{}to paint\textquotedblright{}, from Old High German \emph{mālōn} \textquotedblleft{}to mark, to paint\textquotedblright{}; EtymDB lists the noun without a source.
\end{cousinweb}

\begin{grammarlens}[title={What you've built}]
One more word for telling, explaining and describing at B1.
\end{grammarlens}

\subsection*{Your turn}
\begin{itemize}
\raggedright
  \item \emph{Say it:} \emph{das Gemälde}
  \item \emph{Say it:} \emph{das Gemälde}, once more
  \item \emph{Say it:} a sentence of your own with \emph{das Gemälde}
  \item \emph{Recall:} say the German for to applaud, then the German for applause, then say \emph{das Gemälde} again
\end{itemize}

\begin{tcolorbox}[breakable,colback=teal!4,colframe=teal!35!black,title={Before you move on}]
What does das Gemälde mean? (\textquotedblleft{}Painting\textquotedblright{}.) Say it once more.
\end{tcolorbox}
\section[der Pinsel]{der Pinsel — paintbrush}

\label{lesson:GE-C518-pinsel}



\begin{quote}
Before the new one: say the German for applause, then the German for painting.
\end{quote}

\subsection*{You'll want to know: der Pinsel}
\textbf{der Pinsel} — \textquotedblleft{}paintbrush\textquotedblright{}.

\begin{quote}
Für die kleinen Details nimmt sie einen dünnen Pinsel.
\end{quote}

\begin{quote}
\emph{For the small details she uses a thin brush.}
\end{quote}

Plural: \emph{die Pinsel}.

\begin{cousinweb}
\emph{Pinsel} comes from Middle High German \emph{bensel}, \emph{pensel}, from Old French \emph{pincel}, from Latin \emph{pēniculus} \textquotedblleft{}little tail, brush\textquotedblright{} (EtymDB).
\end{cousinweb}

\begin{grammarlens}[title={What you've built}]
One more word for telling, explaining and describing at B1.
\end{grammarlens}

\subsection*{Your turn}
\begin{itemize}
\raggedright
  \item \emph{Say it:} \emph{der Pinsel}
  \item \emph{Say it:} \emph{der Pinsel}, once more
  \item \emph{Say it:} a sentence of your own with \emph{der Pinsel}
  \item \emph{Recall:} say the German for applause, then the German for painting, then say \emph{der Pinsel} again
\end{itemize}

\begin{tcolorbox}[breakable,colback=teal!4,colframe=teal!35!black,title={Before you move on}]
What does der Pinsel mean? (\textquotedblleft{}Paintbrush\textquotedblright{}.) Say it once more.
\end{tcolorbox}
\section[die Leinwand]{die Leinwand — canvas; screen}

\label{lesson:GE-C518-leinwand}



\begin{quote}
Before the new one: say the German for painting, then the German for paintbrush.
\end{quote}

\subsection*{You'll want to know: die Leinwand}
\textbf{die Leinwand} — \textquotedblleft{}canvas; screen\textquotedblright{}.

\begin{quote}
Der Maler stellt eine neue Leinwand auf.
\end{quote}

\begin{quote}
\emph{The painter sets up a new canvas.}
\end{quote}

Plural: \emph{die Leinwände}. Also a cinema screen.

\begin{cousinweb}
\emph{Leinwand} is inherited from Middle High German \emph{līnwāt} \textquotedblleft{}linen\textquotedblright{} (EtymDB), built on \emph{lein} \textquotedblleft{}flax\textquotedblright{} and \emph{wāt} \textquotedblleft{}cloth\textquotedblright{}; the second part was later reshaped like \emph{Wand} \textquotedblleft{}wall\textquotedblright{}.
\end{cousinweb}

\begin{grammarlens}[title={What you've built}]
One more word for telling, explaining and describing at B1.
\end{grammarlens}

\subsection*{Your turn}
\begin{itemize}
\raggedright
  \item \emph{Say it:} \emph{die Leinwand}
  \item \emph{Say it:} \emph{die Leinwand}, once more
  \item \emph{Say it:} a sentence of your own with \emph{die Leinwand}
  \item \emph{Recall:} say the German for painting, then the German for paintbrush, then say \emph{die Leinwand} again
\end{itemize}

\begin{tcolorbox}[breakable,colback=teal!4,colframe=teal!35!black,title={Before you move on}]
What does die Leinwand mean? (\textquotedblleft{}Canvas; screen\textquotedblright{}.) Say it once more.
\end{tcolorbox}
\section[die Staffelei]{die Staffelei — easel}

\label{lesson:GE-C518-staffelei}



\begin{quote}
Before the new one: say the German for paintbrush, then the German for canvas.
\end{quote}

\subsection*{You'll want to know: die Staffelei}
\textbf{die Staffelei} — \textquotedblleft{}easel\textquotedblright{}.

\begin{quote}
Die Staffelei steht am Fenster, wo das Licht gut ist.
\end{quote}

\begin{quote}
\emph{The easel stands by the window, where the light is good.}
\end{quote}

Plural: \emph{die Staffeleien}.

\begin{cousinweb}
\emph{Staffelei} is built on \emph{Staffel} \textquotedblleft{}step, rung\textquotedblright{}, which EtymDB traces to Proto-Germanic \emph{stapulaz} \textquotedblleft{}post\textquotedblright{}; the ending \emph{-ei} makes a noun.
\end{cousinweb}

\begin{grammarlens}[title={What you've built}]
One more word for telling, explaining and describing at B1.
\end{grammarlens}

\subsection*{Your turn}
\begin{itemize}
\raggedright
  \item \emph{Say it:} \emph{die Staffelei}
  \item \emph{Say it:} \emph{die Staffelei}, once more
  \item \emph{Say it:} a sentence of your own with \emph{die Staffelei}
  \item \emph{Recall:} say the German for paintbrush, then the German for canvas, then say \emph{die Staffelei} again
\end{itemize}

\begin{tcolorbox}[breakable,colback=teal!4,colframe=teal!35!black,title={Before you move on}]
What does die Staffelei mean? (\textquotedblleft{}Easel\textquotedblright{}.) Say it once more.
\end{tcolorbox}
\section[das Aquarell]{das Aquarell — watercolour}

\label{lesson:GE-C518-aquarell}



\begin{quote}
Before the new one: say the German for canvas, then the German for easel.
\end{quote}

\subsection*{You'll want to know: das Aquarell}
\textbf{das Aquarell} — \textquotedblleft{}watercolour\textquotedblright{}.

\begin{quote}
Im Urlaub malt mein Vater Aquarelle von der Küste.
\end{quote}

\begin{quote}
\emph{On holiday my father paints watercolours of the coast.}
\end{quote}

Plural: \emph{die Aquarelle}.

\begin{cousinweb}
\emph{Aquarell} came from French \emph{aquarelle} (EtymDB), from Italian \emph{acquarella}, built on Latin \emph{aqua} \textquotedblleft{}water\textquotedblright{}.
\end{cousinweb}

\begin{grammarlens}[title={What you've built}]
One more word for telling, explaining and describing at B1.
\end{grammarlens}

\subsection*{Your turn}
\begin{itemize}
\raggedright
  \item \emph{Say it:} \emph{das Aquarell}
  \item \emph{Say it:} \emph{das Aquarell}, once more
  \item \emph{Say it:} a sentence of your own with \emph{das Aquarell}
  \item \emph{Recall:} say the German for canvas, then the German for easel, then say \emph{das Aquarell} again
\end{itemize}

\begin{tcolorbox}[breakable,colback=teal!4,colframe=teal!35!black,title={Before you move on}]
What does das Aquarell mean? (\textquotedblleft{}Watercolour\textquotedblright{}.) Say it once more.
\end{tcolorbox}
